\section[Application of the existence theorem for sheaves]{Application of the existence theorem for sheaves:
semicontinuity theorem for the fundamental groups
of the fibers of a proper and separable morphism}
\label{X.2}
% original p. 268

\begin{theorem}
\label{X.2.1}
Let $Y$ be the spectrum of a \emph{complete} noetherian local ring,
with residue field~$\kres$, $X$ a proper $Y$-scheme,
$X_0=X\otimes_A \kres$, $a_0$ a geometric point of
$X_0$ and $a$ the corresponding geometric point of~$X$. Then
the canonical homomorphism $\pi_1(X_0,a_0)\to \pi_1(X,a)$ is an
\emph{isomorphism.}
\end{theorem}

This is only a translation, into the language of the fundamental group,
of the result recalled in IX~\Ref{IX.1.10}. It is here that the
existence theorem for sheaves in formal algebraic geometry enters in
an essential way into the theory of the fundamental group.

Let us now introduce an algebraic closure $\bar{k}$ of the
residue field $k$, and the geometric fiber $\overline X_0
= X_0\otimes_k \bar{k}$. We thus have the exact sequence (IX~\Ref{IX.6.1})
$$
e\to \pi_1(\overline X_0,\overline a)\to \pi_1(X_0,a_0)\to
\pi_1(k,\bar{k})\to e,
$$
on the other hand we have the isomorphism~\Ref{X.2.1} and the
analogous, more elementary, isomorphism
$$
\pi_1(k,\bar{k})\to \pi_1(Y,b),
$$
where $b$ is the image of~$a$ in~$Y$. We thus find:

\begin{corollary}
\label{X.2.2}
With the preceding notation, suppose
$\overline X_0=X_0\otimes_k \bar{k}$
connected, and let $\overline a_0$ be a geometric point of
$\overline X_0$, $a_0$ its image in~$X$,
$b_0$ its image in~$Y$. Then the following sequence of canonical
homomorphisms
$$
e\to \pi_1(\overline X_0,\overline a)\to \pi_1(X,a_0)\to
\pi_1(Y,b_0)\to e
$$
is exact.
\end{corollary}

One will compare this sequence with the exact sequence~\Ref{X.1.4}, but
one will note that~a)~we have not had to make any hypothesis of flatness,
or of separability on the fibers, for $X\to Y$; b)~we have the
important complement that \emph{the morphism $\pi_1(\overline
X_0,\overline a_0)\to \pi_1(X,a_0)$ is injective.}

This last fact will allow us to compare the fundamental group of the
other
% original p. 269
geometric fibers of~$X$ over~$Y$ with that of~$\overline
X_0$. Indeed, let $y_1$ be an arbitrary point of~$Y$, $X_1$ its fiber
and $\overline X_1$ its geometric fiber, relative to
an algebraically closed extension of~$\kres(y_1)$, $\overline a_1$
a geometric point of~$\overline X_1$, $a_1$ its image in
$X$ and $b_1$ its image in~$Y$. Let us choose a ``class of paths''
from~$a_1$ to $a_0$, whence a class of paths from~$b_1$ to
$b_0$, whence a commutative diagram of homomorphisms:
$$
\xymatrix@=.5cm{
&\pi_1(\overline X_1,\overline a_1) \ar[r] & \pi_1(X,a_1) \ar[r]\ar[d] & \pi_1(Y,b_1) \ar[r]\ar[d] & e \\
e \ar[r]& \pi_1(\overline X_0,\overline a_0) \ar[r] & \pi_1(X,a_0) \ar[r] &\pi_1(Y,b_0) \ar[r] & e,
}
$$
where the two vertical arrows drawn are
isomorphisms. Since the second row is exact, we therefore find
a canonical homomorphism, which we shall call \emph{the specialization
homomorphism for the fundamental group} (depending only
on the class of paths chosen from~$a_1$ to~$a_0$, hence
\emph{defined modulo inner automorphism of}
$\pi_1(X,a_0)$):
$$
\pi_1(\overline X_1,\overline a_1)\to \pi_1(\overline X_0,\overline
a_0).
$$
When the first row above is also exact, it
follows at once that the specialization homomorphism
is surjective. We therefore find, taking~\Ref{X.1.4} into account:

\begin{corollary}
\label{X.2.3}
Under the conditions of~\Ref{X.2.1}, suppose moreover that the morphism
$f\colon X\to Y$ is \emph{separable} \eqref{X.1.1} and $\overline
X_0$ connected (so that by virtue of~\Ref{X.1.2} we have
$f_*(\cal{O}_X)=\cal{O}_Y$). Then for every geometric
fiber $\overline X_1$ of~$X$ over~$Y$, equipped with a geometric
point $\overline a_1$, the specialization homomorphism
defined above is a \emph{surjective} homomorphism.
\end{corollary}
This is a \emph{semicontinuity} result for the
fundamental group, which does not yet seem to have an analogue in
algebraic topology. One can moreover state it under more
general conditions:

\begin{corollary}
\label{X.2.4}
Let
% original p. 270
$f\colon X\to Y$ be a proper morphism with geometrically connected
fibers, with $Y$ locally noetherian,
$y_0$ and $y_1$ two points of~$Y$ such that $y_0\in\overline{\{y_1\}}$,
$\overline X_0$ and $\overline X_1$ the geometric fibers of
$X$ corresponding to given algebraically closed extensions
of~$\kres(y_0)$ and $\kres(y_1)$, $\overline a_0$
\resp $\overline a_1$ a geometric point of~$\overline X_0$
\resp $\overline X_1$. Then one can define in a
natural way a specialization homomorphism:
$$
\pi_1(\overline X_1,\overline a_1)\to \pi_1(\overline X_0,\overline
a_0),
$$
defined up to inner automorphism, and this is
a \emph{surjective} homomorphism if $f$ is a separable
morphism \eqref{X.1.1}.
\end{corollary}

Indeed, it follows first of all from~\Ref{X.1.8} that \Ref{X.2.4} is
essentially independent of the algebraically closed extensions
chosen for the residue fields $\kres(y_0)$
and~$\kres(y_1)$. This allows us to replace $Y$ by a
prescheme $Y'$ over~$Y$ having a point $y_0'$ (\resp~$y_1'$)
above~$y_0$ (\resp $y_1$). We shall then take for $Y'$ the
spectrum of the completion of the local ring of~$y_0$ in~$Y$, and we
apply~\Ref{X.2.3}.

\begin{remarks}
\label{X.2.5}
The final conclusion of~\Ref{X.2.4} on the surjectivity of
the specialization homomorphism, and a fortiori the
results~\Ref{X.1.3} and~\Ref{X.1.4} of which it is a
consequence, becomes incorrect if one no longer assumes that $f\colon
X\to Y$ is a separable morphism, even for
projective schemes over an algebraically closed field of
characteristic~$0$. We shall see examples of this later, both
in the case where $f$ is flat but where $f$ admits a non-separable
fiber ($X$ and $Y$ being nevertheless smooth over~$k$), and
in the case where the fibers of~$f$ are indeed separable but
where $f$ is not flat (for example $f\colon X\to Y$ being a
birational morphism of integral normal schemes),
\cf XI~\Ref{XI.3}. In these examples, it may happen that the
fundamental group of the generic geometric fiber is
zero, but not that of a suitable special geometric
fiber. On the other hand, even if $f\colon X\to Y$ is a
proper separable morphism as in~\Ref{X.2.4}, it
commonly happens that the specialization morphism is
% original p. 271
not an isomorphism. Thus, it is easy to give examples
where $\overline X_1$ is a nonsingular elliptic curve,
(hence its fundamental group is commutative, and its
$\ell$-primary component for a prime number $\ell$ prime to the
characteristic is isomorphic to $\ZZ_\ell^2$, \cf \Ref{XI}),
while $\overline X_0$ consists either of two nonsingular
rational curves meeting in two points, or of
two rational curves tangent at one point, or finally of a
rational curve having a singular point which is a cusp.
(For the complete classification of degenerate elliptic
curves, see the recent work of Kodaira~\cite{X.1} and N\'eron).
One then sees that in these cases, the
fundamental group of~$\overline X_0$ is respectively $\hat\ZZ$,
$e$, $e$, hence ``strictly smaller'' than that of~$\overline
X_1$. We shall however see later, when $f$ is a
\emph{smooth} morphism, a bound for the kernel of the
specialization homomorphism, which implies in particular that if $\kres(y_0)$ is
of \emph{characteristic} 0, the specialization homomorphism
is an isomorphism. But even for a smooth morphism, if the
characteristic of~$\kres(y_0)$ is $>0$, it may happen that
the specialization homomorphism is not an isomorphism,
as one sees for example in the case where $X$ is an abelian
scheme over~$Y$ (of relative dimension~$1$, if one wishes),
\cf XI~\Ref{XI.2}.
A satisfactory theory of the
specialization of the fundamental group must take into account the
``continuous component'' of the ``true'' fundamental group, corresponding
to the classification of principal coverings with
infinitesimal groups as structure group; granting which one would be
entitled to expect that the ``true'' fundamental groups of the
geometric fibers of a smooth and proper morphism $f\colon
X\to Y$ form a nice local system on~$X$, projective limit
of finite and flat group schemes over~$X$\kern1pt\footnote{This
extremely attractive conjecture is unfortunately contradicted
by an unpublished example of M\ptbl Artin, already when
the fibers of~$f$ are algebraic curves of given genus
$g\ge 2$.}. We shall return later to this point of view,
our present object being on the contrary to push as far
as possible the phenomena common to the topological
theory and the scheme-theoretic theory of the fundamental group.
\end{remarks}

Let now $X_0$ be a proper, smooth and connected curve of genus $g$
over an algebraically closed field~$k$. If $k$ is of
characteristic zero, its fundamental group can be
determined by transcendental means in the following way. We
know that $X_0$ comes by extension of the base from a curve
defined
% original p. 272
over an algebraically closed extension of finite
transcendence degree of the prime field $\QQ$, and taking
\Ref{X.1.8} into account, we may assume that $k$ is itself of finite
transcendence degree over~$\QQ$. We may therefore assume that $k$ is a
subfield of the field $\CC$ of complex numbers, and a new
application of~\Ref{X.1.8} allows us to assume that $k=\CC$. It
is then not difficult to verify that the fundamental group of
$X$ is isomorphic to the compactification of the fundamental group of the
associated topological space $\tilde X$ (compact oriented surface of
genus~$g$), for the topology defined by the subgroups
of finite index\footnote{This deduction was made explicit
in one of the oral expos\'es which were not
written up.}. It is on the other hand classical that the
topological fundamental group is generated by $2g$ generators
$s_i,t_i$ ($1\le i\le g$), subject to a single relation:
$$
(s_1t_1s_1^{-1}t_1^{-1})\cdots(s_gt_gs_g^{-1}t_g^{-1})=1.
$$
Hence the fundamental group of~$X$ admits $2g$ \emph{topological}
generators $s_i,t_i$ ($1\le i\le g$), bound by the single
preceding relation. If now the characteristic of
$k$ is $p>0$, let us denote by $A$ the ring of Witt vectors
constructed with~$k$, by $K$ an algebraically closed extension of
its field of fractions. We have seen in III~\Ref{III.7.4} that there
exists a scheme $X$ over~$Y=\Spec(A)$, proper and smooth over~$Y$,
reducing to $X_0$. Applying~\Ref{X.2.3} to it, we find a
\emph{surjective} morphism
$$
\pi_1(X_1)\to \pi_1(X_0),
$$
where $X_1=X\otimes_A K$. It is immediate\footnote{\cf EGA
IV~12.2.} that $X_1$ is smooth over~$K$, connected \eqref{X.1.2}, of
dimension~$1$, and its genus is equal to~$g$ (by
the invariance of the Euler-Poincar\'e characteristic, \cf EGA
III~7). Since $K$ is of characteristic~$0$, one can
apply the preceding result to it. We have thus proved by
\emph{transcendental means}:

\begin{theorem}
\label{X.2.6}
Let $X_0$ be a smooth proper and connected algebraic curve over an
% original p. 273
algebraically closed field $k$, and let $g$ be its genus. Then
$\pi_1(X_0)$ admits a system of~$2g$ topological
generators, bound by the relation written above. When the
characteristic of~$k$ is~$0$, $\pi_1(X_0)$ is even the
free group of Galois type for the generators and the
relation which precede.
\end{theorem}

\begin{remarks}
\label{X.2.7}
There does not exist at the present time, to the knowledge of the
writer, a proof by purely
algebraic means of the preceding result, (except for the genera
$0,1$). To begin with, one hardly sees how to distinguish in
$\pi_1(X_0)$ $2g$~elements, which one could then expect
to form a system of topological
generators. In this respect, the situation of the rational
line minus~$n$ points, and the study of the
coverings of it tamely ramified at these
points, is more congenial, since the consideration of the ramification groups
at these $n$ points furnishes $n$ elements of the
fundamental group to be studied, which one indeed shows
topologically generate this fundamental group, as we shall see
later\footnote{\Cf Exp\ptbl \Ref{XII}. Even so, these
elements are really determined only modulo
conjugation, and it is appropriate to make a judicious \emph{simultaneous}
choice of these elements in their classes.}. But even
in this particularly concrete case, there does not seem to exist
a purely algebraic proof. Such a
proof would obviously be extremely
interesting. The only fact concerning the fundamental group of a
curve that one knows how to prove by purely algebraic means
(with the exception of the weak finiteness theorem~\Ref{X.2.12}
below, proved by algebraic means by Lang-Serre~\cite{X.2}),
seems to be the determination of the abelianized fundamental group
via the Jacobian (pointed out in IX~\Ref{IX.5.8}, last
line).
\end{remarks}

\subsection{}
\label{X.2.8}
The last assertion of~\Ref{X.2.6} is no longer valid in
characteristic $p>0$, as one sees already in the case of
elliptic curves. As we have already pointed out, we
do not know whether the fundamental group of~$X_0$ is topologically of
finite presentation; this seems quite improbable.

\begin{theorem}
\label{X.2.9}
Let $k$ be an algebraically closed field, and $X$ a proper and
connected scheme over~$k$. Then the fundamental group of~$X$ is
topologically finitely generated.
\end{theorem}

We
% original p. 274
proceed by induction on~$n=\dim X$, the assertion
being trivial for~$n\le 0$. Suppose therefore~$n>0$, and the
theorem proved for the
dimensions~$n'<n$. By Chow's lemma (EGA~II~5.6.2) there
exists a projective scheme $X'$ over~$k$ and a surjective morphism
$X'\to X$. We may obviously assume $X'$ reduced, and by
passing to the normalization, normal. Thanks to the theory of
descent, it suffices to prove that the fundamental groups of the
connected components of~$X'$ are topologically finitely
generated (IX~\Ref{IX.5.2}). This therefore brings us back
to the case where $X$ is \emph{projective} and \emph{normal}. If
then $n=1$, it suffices to apply~\Ref{X.2.6}. If $n\ge 2$, we
consider a projective immersion $X\to \PP_k^r$, and a hyperplane
section $Y=X\cdot H$
(endowed
with the induced reduced structure), such that
$Y\ne X$,
\ie $H\not\supset X$. We shall then have $\dim Y<n$, and taking into account
the induction hypothesis, it suffices to prove that
$\pi_1(Y)\to \pi_1(X)$ is \emph{surjective.} Now more
generally:

\begin{lemma}
\label{X.2.10}
Let $X$ be a prescheme proper over an algebraically closed
field~$k$, $g\colon X\to \PP_k^r$ a morphism. Suppose $X$
irreducible and normal and $\dim g(X)\ge 2$. Let $H$ be a
hyperplane of~$\PP_k^r$ and $Y=X\times_{\PP_k^r} H$. Then $Y$ is
connected, and the homomorphism $\pi_1(Y)\to\pi_1(X)$ is surjective.
\end{lemma}

These assertions indeed follow from the following one:

\begin{corollary}
\label{X.2.11}
Under the preceding conditions, let $X'$ be a connected \'etale
covering of~$X$, and $Y'=X'\times_X Y=X'\times_{\PP_k^r} H$ the
induced covering on~$Y$. Then $Y'$ is connected.
\end{corollary}

Since $X$ is normal, $X'$ is normal, hence, being connected, $X'$
is irreducible; moreover its image in $\PP_k^r$ is of
dimension~$\ge 2$. A well-known lemma due to Zariski (and
called ``\emph{Bertini's theorem}'') therefore implies that
if $H_1'$ is the generic hyperplane in~$\PP_k^r$,
defined over an extension $K$ of~$k$, then $X'\times_{\PP^r} H_1$
is universally irreducible hence universally connected
over~$K$. Zariski's connectedness theorem (EGA~III~4)
then implies that for \emph{every} hyperplane $H$ (defined over an
arbitrary extension of~$k$) $X'\times_{\PP^r} H$ is
geometrically connected. This completes the proof
of~\Ref{X.2.11}, hence of~\Ref{X.2.9}.

\begin{corollary}[Lang-Serre]
\label{X.2.12}
Under
% original p. 275
the conditions of~\Ref{X.2.9}, for every
finite group~$G$, the set of classes, up to isomorphism,
of principal coverings of~$X$ with group~$G$, is finite.
\end{corollary}

\begin{remark}
\label{X.2.13}
Under the conditions of~\Ref{X.2.10} we shall prove, when $\dim
g(X)\ge 3$ (at least when $g$ is an immersion and $X$
regular), a more precise result, known in
algebraic geometry under the name of
``\emph{Lefschetz theorem}'': $\pi_1(Y)\to \pi_1(X)$
\emph{is an isomorphism.}\footnote{\Cf the seminar SGA~2 (1962)
following this one, theorem X~3.10.}
There are, in the classical cases, analogous
statements for the homology groups and the higher
homotopy groups, which sooner or later will have to be
encompassed in abstract algebraic
geometry. Even for the Hodge cohomology $\H^p(X,\Omega^q)$,
it does not seem that the question has yet been
studied; it is moreover hardly probable that for this
latter, the Lefschetz theorems remain valid as they stand
in characteristic~$p>0$.
\end{remark}

\begin{remarkMR}
\label{X.2.14}
Let $R$ be a complete discrete valuation ring, with
algebraically closed residue field $\kres$ of characteristic $p>0$, with
field of fractions~$K$, and let $Y$ be a proper and smooth,
connected curve of genus $g$ over~$R$. We thus have a surjective
specialization morphism $\sp : \pi_1(Y_{\overline{K}}) 
\to\pi_1(Y_{\kres})$.
We have already noted that if $K$ is of characteristic 0, $\sp$ is not an
isomorphism as soon as $g\geq1$. Suppose~$K$ of characteristic $p$,
so that $R$ is isomorphic to the ring of formal
power series~$\kres[[T]]$.
In equal characteristic $p >0$, the fundamental group is not
determined by the genus~$g$, as one already sees on elliptic curves,
which may be ordinary or supersingular. Let us cite
the recent result of A\ptbl Tamagawa (not yet published). If~$G$ is a
profinite group, we denote by $G^{\res}$ the profinite quotient group of~$G$,
projective limit of the topological quotients of~$G$ which are finite
solvable.

\begin{theoremstar}[A\ptbl Tamagawa]
Suppose $g\geq2$, that the special fiber $Y_{\kres}$
is definable over a finite field, and that the morphism
$\sp^{\res}:\pi_1(Y_{\overline{K}})^{\res}\to\pi_1(Y_{\kres})^{\res}$ deduced from
$\sp$ by passage to the quotient is bijective. Then the curve $Y$ is
constant over~$R$.
\end{theoremstar}

Let us note that the Galois group of~$\overline{K}/K$ is solvable. The
preceding statement can also be translated as follows: suppose that
every finite \'etale covering $X_K\to Y_K$ of the generic fiber $Y_K$,
Galois with solvable Galois group, has potentially good
reduction (that is, extends to a finite \'etale covering of~$Y$
after possible replacement of~$R$ by its normalization in a
finite extension of~$K$). Then $Y$ is constant over~$R$.
\end{remarkMR}
